%========================
% chapters/01_introduction.tex
%========================
\chapter{Introduction}
\label{chap:introduction}

Graphs are a natural representation for systems whose meaning is defined not only by individual entities, but also by the pattern of relations among them. Molecules, social networks, communication systems, power grids, circuits, knowledge graphs, and biological structures are all examples in which topology is part of the object being modeled. This makes graph generation a central task in modern machine learning: a generator should not merely reproduce attributes attached to vertices or edges, but should also learn how local decisions combine into global connectivity, degree distribution, motifs, and structural validity.

Generative models for graphs face a difficulty that is less visible in image or text generation. A graph is invariant to node relabeling, whereas a neural model usually receives a tensor, a sequence, or another ordered representation. When an unordered graph is converted into an ordered training object, the conversion is not unique. The same graph with $n$ nodes admits up to $n!$ node permutations, and each permutation can induce a different sequence of construction actions. Consequently, a graph generator that appears to model a graph distribution may in fact be modeling a distribution over graph serializations. This tension is especially important in autoregressive graph generation, where the model builds a graph step by step and every decision is conditioned on the partial graph already produced.

This dissertation studies that tension through the lens of node ordering. The central argument is that node ordering is not a harmless implementation detail. It determines which vertices are visible when an edge must be predicted, which candidate destinations compete in the model's output distribution, how long-range dependencies appear in the action sequence, and how regular the stop decisions become. In a model such as the Deep Generative Model of Graphs (DGMG) \cite{li2018learning}, the generator alternates among decisions to add a node, add an edge, choose an edge destination, and stop. The ordering of the original graph determines the supervised action sequence used during training. Two orderings of the same graph therefore expose the same topology through different conditional histories, and these histories can vary substantially in difficulty.

The work begins with a controlled reproduction of the graph-to-sequence mechanism and uses this reproduction to compare deterministic orderings across several graph families. Degree ordering, breadth-first search, depth-first search, degeneracy ordering, LexBFS, maximum cardinality search, spectral ordering, and random permutations are evaluated on cycles, trees, binary trees, stars, and Barab\'asi--Albert graphs. This stage shows that different graph families favor different orderings. Cycles benefit from orderings that nearly linearize the ring; trees and binary trees favor level-wise parent-child locality; stars benefit from exposing the hub early; and Barab\'asi--Albert graphs reveal trade-offs among size control, connectivity, and hub structure. These results motivate a stronger claim: because no single handcrafted ordering is uniformly best, ordering should itself become a learnable object.

The main methodological contribution is therefore a deep reinforcement-learning formulation for node ordering. Instead of imposing a fixed ordering rule, a policy observes the input graph and selects one node at a time until a complete permutation is produced. The resulting ordering is translated into a graph generation action sequence, and the quality of this sequence is assessed through downstream generator performance. Since the choice of a node is discrete and the reward is only available after the generator has been trained and evaluated, direct differentiation through the ordering is not possible. Reinforcement learning provides a natural solution: the ordering policy is treated as a stochastic graph-conditioned policy and optimized with a policy-gradient estimator.

The learned-ordering method used in the final experiments is called Attention RL ordering. Its policy network encodes the graph with message passing and attention, maintains the set of already selected nodes, masks invalid actions, and produces a probability distribution over the remaining nodes. The policy is trained in an alternating loop with the autoregressive generator, which in our case of study is the DGMG \cite{li2018learning}, but it could be any other generator; The choice of using DGMG was merely because it is one of the foundational approaches to graph generation through decision sequences. At each outer iteration, the policy samples orderings for the training graphs, the translator produces action sequences, the generator is updated on those sequences, and validation and sampling metrics are used to update the policy. The final version strengthens the earlier controller with three graph layers, an attention mechanism, and a degree-based warm start.

The empirical focus of the final learned-ordering study is the Barab\'asi--Albert family \cite{barabasi1999emergence}. This family is a useful test bed because it combines growth, preferential attachment, hub dominance, sparse connectivity, and variable graph size. Unlike cycles or trees, it does not admit a single exact validity predicate that fully captures structural quality. The evaluation therefore uses several complementary metrics: average generated size, valid-size ratio, connected ratio, degree Gini coefficient, hubiness, and average number of edges. The final Attention RL ordering run achieves perfect valid-size ratio, near-perfect connectivity, and realistic size, while keeping degree inequality and hub-related statistics competitive with the deterministic baselines. This supports the view that a learned ordering policy can act as a flexible structural prior for sequential graph generation.

The dissertation is organized as follows. Chapter~\ref{chap:foundations} introduces the graph-theoretic, generative-modeling, and reinforcement-learning concepts used throughout the work. Chapter~\ref{chap:related-work} positions the dissertation with respect to graph generative models, autoregressive graph generation, ordering-aware methods, and reinforcement learning. Chapter~\ref{chap:dgmg-reproduction} describes the DGMG reproduction and the graph-to-sequence translator that makes the ordering problem explicit. Chapter~\ref{chap:deterministic-orderings} analyzes deterministic orderings and explains how they reshape the action-sequence learning problem. Chapter~\ref{chap:rl-ordering} presents the learned-ordering formulation, including the Markov Decision Process, the policy architecture, the reward, the policy-gradient theorem, and the alternating training algorithm. Chapter~\ref{chap:experimental-protocol} defines the datasets, baselines, metrics, and implementation protocol. Chapter~\ref{chap:results-discussion} reports and interprets the deterministic and learned-ordering results. Finally, Chapter~\ref{chap:conclusion} summarizes the contributions, limitations, and future research directions.
